\documentclass[11pt]{article}

\usepackage[spanish]{babel}
\usepackage[utf8]{inputenc}
\usepackage[T1]{fontenc}
\usepackage{geometry}
\usepackage{multicol}
\usepackage{setspace}
\geometry{margin=2cm}

\title{Informe: Metodologías Ágiles, Scrum y Niveles Belt}
\author{}
\date{}

\begin{document}

\maketitle

\begin{multicols}{2}

\section*{Introducción}
En la actualidad, los proyectos requieren adaptarse rápidamente a cambios y nuevas necesidades. 
Por ello surgen las metodologías ágiles, que permiten desarrollar productos de forma flexible, 
iterativa y colaborativa, mejorando la eficiencia y reduciendo errores.

\section*{Metodologías Ágiles}
Las metodologías ágiles son enfoques de trabajo que priorizan la entrega continua de valor 
y la adaptación constante.

\textbf{Características principales:}
\begin{itemize}
    \item Trabajo en ciclos cortos
    \item Entregas frecuentes
    \item Adaptación a cambios
    \item Comunicación constante con el cliente
    \item Mejora continua
\end{itemize}

\textbf{Ventajas:}
\begin{itemize}
    \item Mayor flexibilidad
    \item Detección temprana de errores
    \item Mayor satisfacción del cliente
    \item Trabajo colaborativo
\end{itemize}

\section*{Scrum}
Scrum es uno de los marcos de trabajo ágiles más utilizados. Se basa en dividir el trabajo 
en iteraciones llamadas \textbf{Sprints}, que suelen durar entre 1 y 4 semanas.

\textbf{Roles:}
\begin{itemize}
    \item \textbf{Product Owner:} define los requisitos del proyecto.
    \item \textbf{Scrum Master:} facilita el proceso y elimina obstáculos.
    \item \textbf{Equipo de desarrollo:} realiza las tareas necesarias.
\end{itemize}

\textbf{Eventos:}
\begin{itemize}
    \item \textbf{Sprint Planning:} planificación del trabajo.
    \item \textbf{Daily Scrum:} reunión diaria breve.
    \item \textbf{Sprint Review:} revisión del resultado.
    \item \textbf{Sprint Retrospective:} análisis y mejora del proceso.
\end{itemize}

\textbf{Artefactos:}
\begin{itemize}
    \item Product Backlog
    \item Sprint Backlog
    \item Incremento del producto
\end{itemize}

\section*{Lean Six Sigma y Niveles Belt}
Lean Six Sigma es una metodología enfocada en la mejora de procesos y la reducción de errores.

\textbf{White Belt:}
\begin{itemize}
    \item Nivel básico
    \item Conocimiento general
    \item Participación en proyectos
\end{itemize}

\textbf{Green Belt:}
\begin{itemize}
    \item Nivel intermedio
    \item Análisis de datos
    \item Mejora de procesos
    \item Liderazgo de proyectos pequeños
\end{itemize}

\section*{Relación entre metodologías}
Las metodologías ágiles y Lean Six Sigma pueden complementarse:
\begin{itemize}
    \item Scrum organiza el trabajo
    \item Lean Six Sigma optimiza los procesos
\end{itemize}

\section*{Conclusión}
Las metodologías ágiles permiten adaptarse a los cambios y mejorar la eficiencia en los proyectos. 
Scrum facilita la organización del trabajo en ciclos cortos, mientras que Lean Six Sigma permite 
mejorar la calidad y reducir errores. Su uso conjunto potencia los resultados.

\end{multicols}

\end{document}